\documentclass[10pt]{article}
\usepackage[T1]{fontenc}\usepackage[utf8]{inputenc}\usepackage{textcomp}
\usepackage[scaled]{helvet}\renewcommand{\familydefault}{\sfdefault}
\DeclareUnicodeCharacter{03B5}{\ensuremath{\varepsilon}}\DeclareUnicodeCharacter{03B2}{\ensuremath{\beta}}
\DeclareUnicodeCharacter{03B7}{\ensuremath{\eta}}\DeclareUnicodeCharacter{03B1}{\ensuremath{\alpha}}
\DeclareUnicodeCharacter{03C7}{\ensuremath{\chi}}\DeclareUnicodeCharacter{2212}{\ensuremath{-}}
\DeclareUnicodeCharacter{2265}{\ensuremath{\geq}}\DeclareUnicodeCharacter{2264}{\ensuremath{\leq}}
\DeclareUnicodeCharacter{2192}{\ensuremath{\rightarrow}}
\usepackage[a4paper,margin=20mm]{geometry}
\usepackage{graphicx}
\pagestyle{empty}\setlength{\parindent}{0pt}\setlength{\parskip}{3pt}
\begin{document}
{\large\bfseries Genome-wide association and haplotype analysis identify candidate genes for grain nutritional quality in wild barley (\textit{Hordeum spontaneum})\par}\vspace{3mm}
Theoretical and Applied Genetics\par
[AUTHOR NAMES]\par
Corresponding author: [CORRESPONDING AUTHOR], [AFFILIATION: institution, department, city, country]; e-mail: [E-MAIL]\par\vspace{10mm}
{\large\bfseries Supplementary Figures\par}\vspace{3mm}
Supplementary Figures S1–S5 of the manuscript, numbered in the order in which they are first cited, each with its legend. The tables (Tables S1–S22) are in the accompanying file of Supplementary Tables.\par\vspace{4mm}
{\small \textbf{Fig. S1} Decay of linkage disequilibrium with physical distance across the genome (\textbf{a}) and per chromosome (\textbf{b})\hfill pages 2–3

\textbf{Fig. S2} crosshap haplotype plots of \textit{GPAT6}, \textit{GH17}, \textit{PHT4;3} and the GDSL esterase/lipase gene\hfill pages 4–8

\textbf{Fig. S3} Geographic origin of the accessions carrying the lead-SNP alleles and the haplotypes of the four candidate genes\hfill page 9

\textbf{Fig. S4} Desert enrichment of the lead-SNP alleles and the haplotypes of the four candidate genes\hfill page 10

\textbf{Fig. S5} Genotype BLUPs of the four nutritional traits across the 29 sampling sites, by ecological region\hfill page 11\par}
\end{document}
